\documentclass[11pt,a4paper]{article}

\usepackage[utf8]{inputenc}
\usepackage[T1]{fontenc}
\usepackage{lmodern}
\usepackage[english]{babel}
\usepackage[margin=2.5cm]{geometry}
\usepackage{estilocrispropioLibro}

\pagestyle{fancy}
\fancyhf{}
\lhead{Physics and Chemistry --- 3\textsuperscript{rd} ESO}
\rhead{\nouppercase{\rightmark}}
\cfoot{\thepage}

\title{\vspace{-1.5cm}\textbf{Written Tests}\\[4pt]\Large Physics and Chemistry --- 3\textsuperscript{rd} ESO}
\author{}
\date{}

\newcommand{\examheader}[1]{%
\begin{center}
\begin{tabular}{ll}
\textbf{Name:} \makebox[6cm]{\dotfill} & \textbf{Date:} \makebox[3cm]{\dotfill} \\
\end{tabular}\\[4pt]
\textit{Time available: 45 minutes.\quad Total: 10 points.}\\[2pt]
\textit{Each question shows its score in brackets. The evaluation criterion (CE) it targets is noted as a comment in the source file.}
\end{center}
}

\begin{document}
\maketitle
\tableofcontents

%===================================================================
\newpage
\section{Unit 1: Scientific Method}
\subsection*{Exam}
\examheader{Unit 1}
\setcounter{preg}{1}

\examexercise[1]{In your own words, describe the six stages of the scientific method, using as an example the observation that a plant placed near a window grows towards the light.}
% CE: CE2.1

\examexercise[1.5]{Convert \qty{72}{\kmph} into \unit{\mps} using conversion factors, showing every step.}
% CE: CE3.2

\examexercise[1.5]{A block of metal has a mass of \qty{2700}{\g} and a volume of \qty{1000}{\cubic\centi\metre}. Calculate its density in \unit{\kilo\g\per\cubic\metre}, using conversion factors.}
% CE: CE1.2

\examexercise[1]{Write the following quantities in scientific notation: (a) \qty{0.000042}{\kilo\g}; (b) \qty{318000000}{\m}.}
% CE: CE3.2

\examexercise[1.5]{The graph below (which you must sketch) shows the volume of a gas increasing steadily as its temperature rises, following a straight line through the origin. (a) Which variable is independent and which is dependent? (b) Is the relationship directly or inversely proportional? Justify your answer using the shape of the graph.}
% CE: CE3.1

\examexercise[1]{Classify the following as fundamental or derived quantities, justifying each answer: electric current, density, amount of substance, energy.}
% CE: CE1.1

\examexercise[1.5]{If \qty{1}{\kg} of hazelnuts costs 12.80 €, calculate, using conversion factors, the price of \qty{1}{\hecto\g} and of \qty{1}{\deca\g} of hazelnuts.}
% CE: CE1.2

\examexercise[1]{A classmate says: "I read online that drinking lemon water cures colds, and it must be true because thousands of people say so on social media." Explain, using what you know about the scientific method, why this reasoning is not scientifically valid and what would be needed to test the claim properly.}
% CE: CE2.1

\subsection*{Solutions}

\begin{solucioneng}
The six stages, applied to the plant example:
\begin{enumerate}
    \item \textbf{Observation}: the plant's stems bend and grow towards the window.
    \item \textbf{Problem statement}: why does the plant grow in that direction? Does light intensity or direction affect growth direction?
    \item \textbf{Hypothesis formulation}: the plant grows towards the light source because it needs light for photosynthesis.
    \item \textbf{Experimentation}: place identical plants in different lighting conditions (near a window, in a dark room, under a lamp from one side) and record growth direction over several days.
    \item \textbf{Analysis of results}: compare growth direction and rate in each group.
    \item \textbf{Formulation of law/theory/model}: if all plants grow towards the strongest light source, this confirms the hypothesis (phototropism); if not, the hypothesis is revised and the process restarts.
\end{enumerate}
\end{solucioneng}

\begin{solucioneng}
\[
72\ \dfrac{\cancel{\unit{\kilo\metre}}}{\cancel{\unit{\hour}}} \cdot \dfrac{1000\ \unit{\metre}}{1\ \cancel{\unit{\kilo\metre}}} \cdot \dfrac{1\ \cancel{\unit{\hour}}}{3600\ \unit{\second}} = \dfrac{72 \cdot 1000}{3600} = \qty{20}{\mps}
\]
\textbf{Result: \qty{20}{\mps}.}
\end{solucioneng}

\begin{solucioneng}
Convert mass and volume to SI units first:
\[
2700\ \cancel{\unit{\g}} \cdot \dfrac{1\ \unit{\kilo\g}}{1000\ \cancel{\unit{\g}}} = \qty{2.7}{\kilo\g}
\]
\[
1000\ \cancel{\unit{\cubic\centi\metre}} \cdot \dfrac{1\ \unit{\cubic\metre}}{\num{e6}\ \cancel{\unit{\cubic\centi\metre}}} = \qty{1e-3}{\cubic\metre}
\]
\[
\rho = \dfrac{m}{V} = \dfrac{2.7\ \unit{\kilo\g}}{\qty{1e-3}{\cubic\metre}} = \qty{2700}{\kilo\g\per\cubic\metre}
\]
\end{solucioneng}

\begin{solucioneng}
\textbf{a)} \( \qty{0.000042}{\kilo\g} = \qty{4.2e-5}{\kilo\g} \) (decimal point moved 5 places right).\\
\textbf{b)} \( \qty{318000000}{\m} = \qty{3.18e8}{\m} \) (decimal point moved 8 places left).
\end{solucioneng}

\begin{solucioneng}
\textbf{a)} Temperature is the independent variable (X-axis, controlled); volume is the dependent variable (Y-axis, measured).\\
\textbf{b)} The relationship is \textbf{directly proportional}: the graph is a straight line passing through the origin, of the form \( y = m\,x \), so when temperature increases, volume also increases at a constant rate.
\end{solucioneng}

\begin{solucioneng}
\begin{itemize}
    \item \textbf{Electric current}: fundamental (it is one of the 7 base SI quantities).
    \item \textbf{Density}: derived (it combines mass and volume, $\rho = m/V$).
    \item \textbf{Amount of substance}: fundamental (base SI quantity, unit mole).
    \item \textbf{Energy}: derived (combination of mass, length and time; unit joule).
\end{itemize}
\end{solucioneng}

\begin{solucioneng}
Using \( \qty{1}{\kg} = \qty{10}{\hecto\g} \) and \( \qty{1}{\kg} = \qty{100}{\deca\g} \):
\[
1\ \cancel{\unit{\hecto\g}} \cdot \dfrac{1\ \cancel{\unit{\kg}}}{10\ \cancel{\unit{\hecto\g}}} \cdot \dfrac{12.80\text{ €}}{1\ \cancel{\unit{\kg}}} = 1.28 \text{ €}
\]
\[
1\ \cancel{\unit{\deca\g}} \cdot \dfrac{1\ \cancel{\unit{\kg}}}{100\ \cancel{\unit{\deca\g}}} \cdot \dfrac{12.80\text{ €}}{1\ \cancel{\unit{\kg}}} = 0.128 \text{ €}
\]
\end{solucioneng}

\begin{solucioneng}
Popularity ("thousands of people say so") is not scientific evidence. A claim is scientifically valid only after it has been tested through the scientific method: a testable hypothesis must be formulated (e.g. "lemon water reduces cold symptoms compared to a control group") and checked with a controlled experiment comparing people who drink lemon water with a control group under the same conditions, ruling out other explanations (placebo effect, natural recovery). Anecdotes shared on social media are not controlled, repeatable, or falsifiable, so they belong to pseudoscience, not science.
\end{solucioneng}

%===================================================================
\newpage
\section{Unit 2: Kinematics}
\subsection*{Exam}
\examheader{Unit 2}
\setcounter{preg}{1}

\examexercise[1]{Define, in your own words, the terms \textit{displacement} and \textit{distance travelled}. If a runner completes one full lap of a circular \qty{400}{\m} track, what is her displacement and what is the distance she has travelled?}
% CE: CE1.1

\examexercise[1.5]{A cyclist travels at a constant speed of \qty{6}{\mps} for \qty{20}{\min}. Calculate the distance travelled, in metres, using SI units.}
% CE: CE1.2

\examexercise[1.5]{A motorbike accelerates from rest to \qty{72}{\kmph} in \qty{6}{\s}. Calculate its average acceleration in \unit{\mpss}.}
% CE: CE1.2

\examexercise[1.5]{A car travelling at \qty{20}{\mps} brakes and comes to a complete stop in \qty{4}{\s}. (a) Calculate its acceleration. (b) Calculate the braking distance.}
% CE: CE1.2

\examexercise[1.5]{The table below shows the position of a trolley at different times.\newline
\begin{center}
% --- tabla antigua ---
% \begin{tabular}{|c|c|c|c|c|}
% \hline
% \textbf{Time (s)} & 0 & 2 & 4 & 6 \\
% \hline
% \textbf{Position (m)} & 0 & 10 & 20 & 30 \\
% \hline
% \end{tabular}
% --- tabla nueva ---
\begin{tabular}{@{}ccccc@{}}
    \toprule
    \textbf{Time (s)} & \textbf{0} & \textbf{2} & \textbf{4} & \textbf{6} \\
    \midrule
    \textbf{Position (m)} & 0 & 10 & 20 & 30 \\
    \bottomrule
\end{tabular}
\end{center}
(a) What type of motion is this? Justify your answer using the data. (b) Calculate the velocity of the trolley.}
% CE: CE3.1

\examexercise[1]{A position-time graph shows a horizontal straight line for an object. What does this tell you about the object's motion? Explain your reasoning.}
% CE: CE1.1

\examexercise[1]{Explain why, when solving a kinematics problem, it is essential to convert all given quantities to SI units before applying any formula. Illustrate with a short example of what could go wrong if this step is skipped.}
% CE: CE3.2

\examexercise[1]{A student says: "If a car's average speed over a journey is \qty{50}{\kmph}, it must have been travelling at exactly \qty{50}{\kmph} the whole time." Is this statement correct? Explain why or why not, referring to the difference between instantaneous and average velocity.}
% CE: CE2.1

\subsection*{Solutions}

\begin{solucioneng}
\textbf{Displacement} is the difference between the final and initial position ($\Delta x = x_f - x_i$), i.e. the straight-line change in position. \textbf{Distance travelled} is the total length of the path followed.

Since the runner returns to her starting point after one lap, $x_f = x_i$, so her \textbf{displacement is zero} (\qty{0}{\m}). However, the \textbf{distance travelled equals the perimeter of the track}, i.e. \qty{400}{\m}.
\end{solucioneng}

\begin{solucioneng}
Convert time to SI: $t = 20\ \text{min} \times \dfrac{60\ \text{s}}{1\ \text{min}} = \qty{1200}{\s}$.

Using $d = v \cdot t$:
\[
d = \qty{6}{\mps} \times \qty{1200}{\s} = \qty{7200}{\m}
\]
\textbf{Result: the cyclist travels \qty{7200}{\m} (7.2 km).}
\end{solucioneng}

\begin{solucioneng}
Convert final velocity to SI units:
\[
v_f = 72\ \dfrac{\text{km}}{\text{h}} \cdot \dfrac{1000\ \text{m}}{1\ \text{km}} \cdot \dfrac{1\ \text{h}}{3600\ \text{s}} = \qty{20}{\mps}
\]
Since it starts from rest, $v_0 = \qty{0}{\mps}$. Applying $a = \dfrac{v_f - v_0}{t}$:
\[
a = \dfrac{20 - 0}{6} = \qty{3.33}{\mpss}
\]
\end{solucioneng}

\begin{solucioneng}
\textbf{a)} $v_0 = \qty{20}{\mps}$, $v_f = \qty{0}{\mps}$, $t = \qty{4}{\s}$.
\[
a = \dfrac{v_f - v_0}{t} = \dfrac{0-20}{4} = \qty{-5}{\mpss}
\]
\textbf{b)} Using $d = v_0 \cdot t + \frac{1}{2} a t^2$:
\[
d = 20 \times 4 + \frac{1}{2} \times (-5) \times 4^2 = 80 - 40 = \qty{40}{\m}
\]
\end{solucioneng}

\begin{solucioneng}
\textbf{a)} Between every pair of consecutive points, the position increases by \qty{10}{\m} every \qty{2}{\s}: the position-time graph would be a straight line, so this is \textbf{Uniform Rectilinear Motion (URM)}, since equal increases in position occur in equal time intervals (constant velocity).

\textbf{b)} $v = \dfrac{\Delta x}{\Delta t} = \dfrac{30-0}{6-0} = \qty{5}{\mps}$.
\end{solucioneng}

\begin{solucioneng}
A horizontal straight line on a position-time graph means the position does not change over time, so the object is \textbf{at rest} (stationary). This is because the slope of a position-time graph represents the velocity, and a horizontal line has zero slope, i.e. $v = 0$.
\end{solucioneng}

\begin{solucioneng}
Kinematic formulas (e.g. $v = x/t$, $a = \Delta v/\Delta t$, $x = v_0 t + \frac{1}{2}at^2$) are only valid when all quantities are expressed in a consistent system of units, normally SI (metres, seconds, m/s, m/s²). If, for example, a distance is left in kilometres and a time in seconds without converting, the calculated velocity would be numerically wrong by a factor of 1000, even though the formula itself was applied correctly. For instance, using \qty{90}{\km} and \qty{3600}{\s} directly without converting km to m would give $v = 0.025$ instead of the correct \qty{25}{\mps}.
\end{solucioneng}

\begin{solucioneng}
The statement is \textbf{not correct}. \qty{50}{\kmph} is the \textbf{average velocity} over the whole journey, calculated as total displacement divided by total time. The \textbf{instantaneous velocity} (the speed shown on a speedometer at a given moment) can be higher or lower than this at different points of the journey --- for example, the car could slow down at junctions and speed up on open road --- as long as the average over the entire trip works out to \qty{50}{\kmph}.
\end{solucioneng}

%===================================================================
\newpage
\section{Unit 3: Dynamic Forces}
\subsection*{Exam}
\examheader{Unit 3}
\setcounter{preg}{1}

\examexercise[1]{Define \textit{force} and name the four elements needed to fully represent a force. In your notebook, sketch a horizontal force of \qty{8}{\newton} pointing to the left, applied at the origin of an XY coordinate system.}
% CE: CE1.1

\examexercise[1.5]{Three friends push a stalled car in the same direction with forces of 150 N, 120 N, and 100 N. Calculate the magnitude and direction of the resultant force.}
% CE: CE1.2

\examexercise[1.5]{Two forces act perpendicularly on an object: one of \qty{9}{\newton} and another of \qty{12}{\newton}. Calculate the magnitude of the resultant force.}
% CE: CE1.2

\examexercise[1.5]{A spring stretches \qty{4}{\centi\metre} when a force of \qty{12}{\newton} is applied. (a) Calculate the spring constant. (b) Calculate the force needed to stretch the same spring by \qty{10}{\centi\metre}.}
% CE: CE1.2

\examexercise[1]{State Newton's First Law (Law of Inertia) in your own words and explain, using this law, why passengers in a car feel pushed backward when the car suddenly accelerates forward.}
% CE: CE1.1

\examexercise[1.5]{A shopping trolley of mass \qty{15}{\kg} is pushed with a force of \qty{30}{\newton} on a horizontal, frictionless floor. Calculate the acceleration of the trolley.}
% CE: CE1.2

\examexercise[1]{Explain, using Newton's Third Law, what happens when a swimmer pushes against the wall of a pool to start swimming. Identify the action force and the reaction force, naming the two bodies involved.}
% CE: CE1.1

\examexercise[1]{A satellite of mass \qty{500}{\kg} orbits a planet. Explain, without necessarily calculating a numerical value, how the gravitational force between the planet and the satellite would change if: (a) the satellite's mass were doubled; (b) the distance to the planet were doubled. Justify your answer using the universal gravitation formula.}
% CE: CE3.1

\subsection*{Solutions}

\begin{solucioneng}
A \textbf{force} is a physical quantity that represents any action capable of modifying the shape of a body or changing its state of rest or motion. The four elements needed to represent a force are: \textbf{point of application}, \textbf{magnitude}, \textbf{direction}, and \textbf{sense}.

\textit{Sketch}: a horizontal arrow of length proportional to 8 N, starting at the origin and pointing along the negative X-axis, labelled $\vec{F} = -8$ N (or "8 N to the left").
\end{solucioneng}

\begin{solucioneng}
Since all three friends push in the same direction, the forces add algebraically:
\[
F_R = F_1 + F_2 + F_3 = 150 + 120 + 100 = \qty{370}{\newton}
\]
\textbf{The resultant force is \qty{370}{\newton}, in the same direction the friends are pushing.}
\end{solucioneng}

\begin{solucioneng}
Since the forces are perpendicular, we apply the Pythagorean theorem:
\[
F_R = \sqrt{F_1^2 + F_2^2} = \sqrt{9^2 + 12^2} = \sqrt{81+144} = \sqrt{225} = \qty{15}{\newton}
\]
\end{solucioneng}

\begin{solucioneng}
\textbf{a)} Using Hooke's Law $F = k \cdot \Delta l$, with $\Delta l = \qty{4}{\centi\metre} = \qty{0.04}{\m}$:
\[
k = \dfrac{F}{\Delta l} = \dfrac{12\ \unit{\newton}}{0.04\ \unit{\m}} = \qty{300}{\newton\per\metre}
\]
\textbf{b)} With $\Delta l = \qty{10}{\centi\metre} = \qty{0.10}{\m}$:
\[
F = k \cdot \Delta l = 300 \times 0.10 = \qty{30}{\newton}
\]
\end{solucioneng}

\begin{solucioneng}
\textbf{Newton's First Law} states that every body remains at rest or in uniform straight-line motion unless acted upon by an external force (inertia). When the car accelerates forward, the passengers' bodies tend, by inertia, to keep their previous state of motion (relative rest), while the car (and the seat) moves forward under them. Because the passenger's body "resists" the change in motion, it appears to be pushed backward relative to the accelerating car, even though no actual backward force is being applied to the passenger --- it is simply inertia keeping the body in its previous state.
\end{solucioneng}

\begin{solucioneng}
Using Newton's Second Law $\Sigma F = m \cdot a$:
\[
a = \dfrac{F}{m} = \dfrac{30\ \unit{\newton}}{15\ \unit{\kg}} = \qty{2}{\mpss}
\]
\end{solucioneng}

\begin{solucioneng}
By \textbf{Newton's Third Law}, when two bodies interact, they exert equal and opposite forces on each other. The \textbf{action force} is the force the swimmer's feet exert on the wall (pushing backward against it). The \textbf{reaction force} is the force the wall exerts on the swimmer's feet, of equal magnitude but opposite direction (pushing the swimmer forward, away from the wall). The two bodies involved are the \textbf{swimmer} and the \textbf{wall of the pool}.
\end{solucioneng}

\begin{solucioneng}
The gravitational force is given by $F_G = G \cdot \dfrac{m_1 \cdot m_2}{r^2}$.

\textbf{a)} If the satellite's mass $m$ is doubled, and everything else stays constant, $F_G$ is directly proportional to $m$, so the gravitational force would also \textbf{double}.

\textbf{b)} If the distance $r$ is doubled, since $F_G$ is inversely proportional to $r^2$, the force would become $\dfrac{1}{2^2} = \dfrac{1}{4}$ of its original value, i.e. it would be \textbf{reduced to a quarter} of the original force.
\end{solucioneng}

%===================================================================
\newpage
\section{Unit 4: Energy}
\subsection*{Exam}
\examheader{Unit 4}
\setcounter{preg}{1}

\examexercise[1]{Define \textbf{energy} and \textbf{work}, and explain in your own words the relationship between these two concepts.}
% CE: CE1.1

\examexercise[1]{A crane exerts a force of \SI{3000}{\newton} to lift a load and does \SI{15000}{\joule} of work in the process. Calculate the height to which the load has been raised.}
% CE: CE1.2

\examexercise[1.5]{A \SI{800}{\kilogram} car is travelling at \SI{72}{\kilo\metre\per\hour}. Calculate its kinetic energy (remember to convert the speed to \unit{\mps} first).}
% CE: CE1.2

\examexercise[1.5]{A \SI{0.6}{\kilogram} ball is at a height of \SI{4}{\metre} and moving with a speed of \SI{6}{\mps}. Calculate its potential energy, its kinetic energy, and its total mechanical energy. (Use $g = \SI{9.8}{\mpss}$.)}
% CE: CE1.2

\examexercise[1.5]{A skateboarder starting from rest slides down a ramp \SI{3.2}{\metre} high with no friction. Using the principle of conservation of mechanical energy, calculate the speed with which they reach the bottom of the ramp.}
% CE: CE1.2

\examexercise[1]{Look at the electricity bill shown in the unit (Figure of the electricity bill). Explain the difference between the \unit{\kilo\watt} of contracted power and the \unit{\kwh} of energy consumed, and say which of the two is a unit of power and which is a unit of energy.}
% CE: CE3.1

\examexercise[1]{A student says: ``This iron is very hot, so it has a lot of heat.'' Explain, using the correct scientific definitions of \textit{heat} and \textit{temperature}, why this sentence is not entirely correct.}
% CE: CE1.1

\examexercise[1]{An electric water heater has a power of \SI{1500}{\watt}. If it is switched on for \SI{4}{\hour} a day, calculate the energy it consumes per day, in joules and in \unit{\kwh}.}
% CE: CE1.2

\examexercise[1]{Explain why wind energy is classified as a renewable energy source, and mention one advantage and one disadvantage of using wind energy to generate electricity in a region like Castilla-La Mancha.}
% CE: CE6.2

\subsection*{Solutions}

\begin{solucioneng}
\textbf{Energy} is the capacity of a body to produce changes in itself or in its surroundings. \textbf{Work} is the transformation produced by a force, that is, we talk about work when a force transmits energy. They are related because work is the mechanism by which energy is transferred from one body to another (or transformed from one form into another): when a force does work on a body, it changes that body's energy.
\end{solucioneng}

\begin{solucioneng}
Using $W = F \cdot d$, we solve for $d$ (here $d$ is the height $h$, since the crane lifts the load):
\[
d = \frac{W}{F} = \frac{15000}{3000} = \qty{5}{\metre}
\]
\textbf{Answer:} The load has been raised \qty{5}{\metre}.
\end{solucioneng}

\begin{solucioneng}
First convert the speed: $v = 72 \text{ km/h} = \dfrac{72 \times 1000}{3600} = \qty{20}{\mps}$
\[
E_k = \frac{1}{2}mv^2 = \frac{1}{2} \times 800 \times 20^2 = \frac{1}{2} \times 800 \times 400 = \qty{160000}{\joule} = \qty{160}{\kilo\joule}
\]
\textbf{Answer:} The kinetic energy is \qty{160}{\kilo\joule}.
\end{solucioneng}

\begin{solucioneng}
Potential energy:
\[
E_p = mgh = 0.6 \times 9.8 \times 4 = \qty{23.52}{\joule}
\]
Kinetic energy:
\[
E_k = \frac{1}{2}mv^2 = \frac{1}{2} \times 0.6 \times 6^2 = 0.3 \times 36 = \qty{10.8}{\joule}
\]
Total mechanical energy:
\[
E_M = E_p + E_k = 23.52 + 10.8 = \qty{34.32}{\joule}
\]
\textbf{Answers:} $E_p = \qty{23.52}{\joule}$, $E_k = \qty{10.8}{\joule}$, $E_M = \qty{34.32}{\joule}$.
\end{solucioneng}

\begin{solucioneng}
By conservation of mechanical energy, all the initial potential energy converts into kinetic energy at the bottom:
\[
mgh = \frac{1}{2}mv^2 \Rightarrow v = \sqrt{2gh}
\]
\[
v = \sqrt{2 \times 9.8 \times 3.2} = \sqrt{62.72} \approx \qty{7.92}{\mps}
\]
\textbf{Answer:} The skateboarder reaches the bottom at approximately \qty{7.92}{\mps}.
\end{solucioneng}

\begin{solucioneng}
The contracted power (\unit{\kilo\watt}) is the maximum rate of energy delivery the household has hired with the electricity company --- it is a unit of \textbf{power}, i.e.\ energy per unit time. The \unit{\kwh} is the amount of \textbf{energy} actually consumed, obtained by multiplying power by the time it is used. In short: \unit{\kilo\watt} tells us ``how fast'' energy can be delivered, while \unit{\kwh} tells us ``how much'' energy has been used.
\end{solucioneng}

\begin{solucioneng}
The sentence confuses \textbf{temperature} with \textbf{heat}. Temperature is a measure of the average kinetic energy of the particles of the iron (how ``agitated'' its molecules are), while heat is energy transferred between two bodies due to a temperature difference, and depends on the mass and type of material, not only on how hot something feels. A hot iron has a high temperature, but the total heat it could transfer also depends on its mass and specific heat capacity, so a small hot object is not necessarily transferring huge amounts of energy.
\end{solucioneng}

\begin{solucioneng}
\[
E = P \cdot t = 1500 \text{ W} \times (4 \times 3600 \text{ s}) = 1500 \times 14400 = \qty{21600000}{\joule} = \qty{21.6e6}{\joule}
\]
Converting to \unit{\kwh}: $E = 1500 \text{ W} \times 4 \text{ h} = \qty{6000}{\watt\hour} = \qty{6}{\kwh}$
\textbf{Answer:} The heater consumes \qty{21.6}{\mega\joule} (\qty{6}{\kwh}) per day.
\end{solucioneng}

\begin{solucioneng}
Wind energy is renewable because its regeneration rate (the wind keeps blowing due to natural atmospheric processes) equals or exceeds the rate at which we use it, so it does not run out with use, unlike fossil fuels. \textbf{Advantage:} it does not emit greenhouse gases during electricity generation and Castilla-La Mancha has abundant wind resources, making it one of the regions in Spain with most installed wind power. \textbf{Disadvantage:} it is an intermittent source (it depends on wind conditions and cannot be guaranteed at all times), and wind farms have a visual and environmental impact on the landscape and on bird life.
\end{solucioneng}

%===================================================================
\newpage
\section{Unit 5: Material Systems}
\subsection*{Exam}
\examheader{Unit 5}
\setcounter{preg}{1}

\examexercise[1]{Classify the following materials into elements, compounds, homogeneous mixtures or heterogeneous mixtures: oxygen gas ($O_2$), a glass of orange juice with pulp, table salt dissolved in water, and pure iron.}
% CE: CE1.1

\examexercise[1.5]{A metal alloy contains \SI{85}{\gram} of copper and \SI{15}{\gram} of zinc in a total mass of \SI{200}{\gram}. Calculate the mass percentage of copper, of zinc, and of any other component present.}
% CE: CE1.2

\examexercise[1.5]{We dissolve \SI{8}{\gram} of potassium chloride and \SI{40}{\gram} of sugar in \SI{350}{\gram} of water. Calculate the mass percentage of each solute in the resulting solution.}
% CE: CE1.2

\examexercise[1.5]{A solution is prepared with \SI{40}{\milli\litre} of vinegar and \SI{160}{\milli\litre} of water. Calculate the percentage by volume of vinegar and of water in the solution.}
% CE: CE1.2

\examexercise[1.5]{We want to prepare \SI{300}{\milli\litre} of a saline solution with a mass concentration of \SI{18}{\gram\per\litre}. What mass of salt is needed? If the density of the solution is \SI{1.01}{\gram\per\cmcubicos}, what is the total mass of the solution?}
% CE: CE1.2

\examexercise[1]{Using the solubility curves given in the unit, explain what happens (at particle level) when we cool a saturated solution of \ch{KNO3} in water from \SI{60}{\celsius} to \SI{20}{\celsius}. Why does a solid appear at the bottom of the container?}
% CE: CE1.1

\examexercise[1]{Explain, in your own words, the difference between the concentration of a solution (mass/volume) and the density of that solution. Why are they often confused, and what units help us tell them apart?}
% CE: CE3.2

\examexercise[1]{A bottle of mineral water lists on its label: sodium, calcium, magnesium and bicarbonate. Identify which is the solvent and which are the solutes in this solution, and explain how you know.}
% CE: CE1.1

\subsection*{Solutions}

\begin{solucioneng}
\begin{itemize}
    \item Oxygen gas ($O_2$): \textbf{element} (only one type of atom).
    \item Orange juice with pulp: \textbf{heterogeneous mixture} (the pulp can be distinguished with the naked eye from the liquid).
    \item Salt dissolved in water: \textbf{homogeneous mixture} (a solution --- the components cannot be distinguished).
    \item Pure iron: \textbf{element} (only one type of atom, Fe).
\end{itemize}
\end{solucioneng}

\begin{solucioneng}
\textbf{Mass of copper and zinc:} $85 + 15 = \qty{100}{\gram}$, but the total mass of the alloy is \qty{200}{\gram}, so there must be another component (or the values given are for a portion). Working with the data as given:
\[
\%Cu = \frac{85}{200} \times 100 = 42.5\%
\]
\[
\%Zn = \frac{15}{200} \times 100 = 7.5\%
\]
Since $42.5\% + 7.5\% = 50\%$, the remaining component(s) make up $100 - 50 = 50\%$ of the alloy.
\textbf{Answer:} 42.5\% Cu, 7.5\% Zn, and 50\% of other component(s).
\end{solucioneng}

\begin{solucioneng}
Total mass of solution: $8 + 40 + 350 = \qty{398}{\gram}$
\[
\%KCl = \frac{8}{398} \times 100 \approx 2.01\%
\]
\[
\%\text{sugar} = \frac{40}{398} \times 100 \approx 10.05\%
\]
\textbf{Answer:} approximately 2.01\% KCl and 10.05\% sugar.
\end{solucioneng}

\begin{solucioneng}
Total volume of solution: $40 + 160 = \qty{200}{\milli\litre}$
\[
\%V_{\text{vinegar}} = \frac{40}{200} \times 100 = 20\%
\]
\[
\%V_{\text{water}} = \frac{160}{200} \times 100 = 80\%
\]
\textbf{Answer:} 20\% vinegar and 80\% water.
\end{solucioneng}

\begin{solucioneng}
Convert the volume to litres: $300 \text{ mL} = 0.3 \text{ L}$.
\[
m_{\text{salt}} = C \cdot V = 18 \text{ g/L} \times 0.3 \text{ L} = \qty{5.4}{\gram}
\]
For the total mass of the solution, we use the density:
\[
m_{\text{solution}} = \rho \cdot V = 1.01 \text{ g/cm}^3 \times 300 \text{ cm}^3 = \qty{303}{\gram}
\]
\textbf{Answer:} \qty{5.4}{\gram} of salt are needed, and the total mass of the solution is \qty{303}{\gram}.
\end{solucioneng}

\begin{solucioneng}
Solubility decreases as temperature decreases (as shown by the downward slope of the \ch{KNO3} solubility curve). At \SI{60}{\celsius} the solution was saturated, meaning it held the maximum possible amount of dissolved \ch{KNO3} at that temperature. When we cool it to \SI{20}{\celsius}, the solubility of \ch{KNO3} in water is much lower, so the water can no longer keep all the solute particles dissolved. The excess solute particles come out of solution and group together, forming visible solid crystals that settle at the bottom of the container.
\end{solucioneng}

\begin{solucioneng}
\textbf{Concentration} (mass concentration, in g/L) tells us how much \textbf{solute} is dissolved per litre of \textbf{solution}, while \textbf{density} tells us the mass of the \textbf{whole solution} per unit of volume, without distinguishing solute from solvent. They are often confused because both relate mass and volume and can have similar-looking units (g/L for concentration versus g/cm$^3$ or g/mL for density, which are numerically different but conceptually similar). The key is checking what mass is being considered: only the solute's mass (concentration) or the total mass of the solution (density), and paying attention to the units used (g/L is typical of concentration, g/cm$^3$ or g/mL is typical of density).
\end{solucioneng}

\begin{solucioneng}
The \textbf{solvent} is water, because it is present in by far the greatest amount and is the substance that dissolves the rest. The \textbf{solutes} are sodium, calcium, magnesium and bicarbonate, since they are present in much smaller amounts, dissolved within the water. We know this because the definition of solvent is the substance in greater amount in the solution, and the definition of solute is the substance(s) present in smaller amount(s).
\end{solucioneng}

%===================================================================
\newpage
\section{Unit 6: Periodic Table and the Atom}
\subsection*{Exam}
\examheader{Unit 6}
\setcounter{preg}{1}

\examexercise[1]{Indicate the number of protons, neutrons and electrons in a neutral atom of bromine, \ch{^{80}_{35}Br}.}
% CE: CE1.2

\examexercise[1]{An atom has a mass number A = 65 and 35 neutrons. Calculate its atomic number and identify how many electrons it has if it is electrically neutral.}
% CE: CE1.2

\examexercise[1.5]{Indicate the number of protons, neutrons and electrons in the following ions: a) \ch{^{24}_{12}Mg^{2+}}; b) \ch{^{32}_{16}S^{2-}}.}
% CE: CE1.2

\examexercise[1]{Explain what an isotope is and, using the example of chlorine-35 and chlorine-37 (both with Z = 17), justify whether these two atoms are isotopes of each other.}
% CE: CE1.1

\examexercise[1.5]{Write the electron configuration of a neutral silicon atom (Z = 14), and indicate its valence shell and the number of valence electrons.}
% CE: CE3.2

\examexercise[1.5]{Write the electron configuration of the fluoride ion \ch{F^{-}} (Z = 9), and explain which noble gas it is isoelectronic with.}
% CE: CE3.2

\examexercise[1]{Locate the following elements in the periodic table and indicate their group and period: potassium (Z = 19) and chlorine (Z = 17). Explain briefly what physical/chemical family each one belongs to.}
% CE: CE1.1

\examexercise[1]{Elements in the same group of the periodic table show similar chemical properties. Explain why, relating your answer to the electron configuration and the valence electrons of the atoms.}
% CE: CE1.1

\subsection*{Solutions}

\begin{solucioneng}
From the notation \ch{^{80}_{35}Br}: $A = 80$, $Z = 35$.
\begin{itemize}
    \item Protons $= Z = 35$
    \item Neutrons $= A - Z = 80 - 35 = 45$
    \item Electrons $= Z = 35$ (neutral atom)
\end{itemize}
\textbf{Answer:} 35 protons, 45 neutrons, 35 electrons.
\end{solucioneng}

\begin{solucioneng}
Since $A = Z + \text{neutrons}$:
\[
Z = A - \text{neutrons} = 65 - 35 = 30
\]
As the atom is electrically neutral, the number of electrons equals the atomic number:
\[
\text{electrons} = Z = 30
\]
\textbf{Answer:} $Z = 30$ (this is zinc, Zn), and it has 30 electrons.
\end{solucioneng}

\begin{solucioneng}
\textbf{a) \ch{^{24}_{12}Mg^{2+}}}
\begin{itemize}
    \item Protons $= Z = 12$
    \item Neutrons $= A - Z = 24 - 12 = 12$
    \item Electrons: cation with charge 2+ has lost 2 electrons $\Rightarrow 12 - 2 = 10$
\end{itemize}
\textbf{Answer a):} 12 protons, 12 neutrons, 10 electrons.

\textbf{b) \ch{^{32}_{16}S^{2-}}}
\begin{itemize}
    \item Protons $= Z = 16$
    \item Neutrons $= A - Z = 32 - 16 = 16$
    \item Electrons: anion with charge 2- has gained 2 electrons $\Rightarrow 16 + 2 = 18$
\end{itemize}
\textbf{Answer b):} 16 protons, 16 neutrons, 18 electrons.
\end{solucioneng}

\begin{solucioneng}
Isotopes are atoms of the same chemical element (same atomic number $Z$) but with a different number of neutrons, and therefore a different mass number $A$. Chlorine-35 (\ch{^{35}_{17}Cl}) and chlorine-37 (\ch{^{37}_{17}Cl}) both have $Z = 17$ (17 protons), so they are the same element, chlorine. However, chlorine-35 has $35 - 17 = 18$ neutrons, while chlorine-37 has $37 - 17 = 20$ neutrons: different numbers of neutrons and different mass numbers. \textbf{Therefore, yes, chlorine-35 and chlorine-37 are isotopes of each other.}
\end{solucioneng}

\begin{solucioneng}
Silicon has $Z = 14$, so it has 14 electrons. Following the Moeller diagram:
\[
1s^2\ 2s^2\ 2p^6\ 3s^2\ 3p^2
\]
Verification: $2+2+6+2+2 = 14$ electrons. \checkmark
\begin{itemize}
    \item The outermost (highest) shell with electrons is $n = 3$, so the \textbf{valence shell is $n = 3$}.
    \item Electrons in that shell: $3s^2 + 3p^2 = 2 + 2 = 4$, so silicon has \textbf{4 valence electrons}.
\end{itemize}
\end{solucioneng}

\begin{solucioneng}
Fluorine has $Z = 9$, so the neutral atom has 9 electrons: $1s^2\ 2s^2\ 2p^5$. As \ch{F^{-}} it has gained 1 electron, giving $9 + 1 = 10$ electrons:
\[
2p^5 + 1e^- \rightarrow 2p^6
\]
\[
\boxed{\ch{F^{-}:} \quad 1s^2\ 2s^2\ 2p^6}
\]
Verification: $2+2+6 = 10$ electrons. \checkmark
This configuration has 10 electrons, the same as neon (Ne, $Z=10$), so \ch{F^{-}} is \textbf{isoelectronic with neon}.
\end{solucioneng}

\begin{solucioneng}
\textbf{Potassium (K, Z = 19):} electron configuration $[\text{Ar}]\ 4s^1$. The outermost shell is $n = 4$ $\rightarrow$ \textbf{Period 4}. It has 1 valence electron $\rightarrow$ \textbf{Group 1}. Potassium belongs to the family of \textbf{alkali metals}, very reactive metals that react violently with water.

\textbf{Chlorine (Cl, Z = 17):} electron configuration $1s^2\ 2s^2\ 2p^6\ 3s^2\ 3p^5$. The outermost shell is $n = 3$ $\rightarrow$ \textbf{Period 3}. It has 7 valence electrons ($3s^2 3p^5$) $\rightarrow$ \textbf{Group 17}. Chlorine belongs to the family of \textbf{halogens}, very reactive nonmetals.
\end{solucioneng}

\begin{solucioneng}
Elements in the same group of the periodic table occupy the same vertical column, which means they have the \textbf{same number of valence electrons} (electrons in their outermost shell), even though they have different numbers of total electrons and different valence shells ($n$). Since chemical properties and the way an atom forms bonds or reacts depend mainly on its valence electrons (they are the electrons involved in chemical reactions), elements with the same number of valence electrons tend to behave chemically in a similar way --- for example, all alkali metals (group 1) have 1 valence electron and all react easily by losing that electron, while all halogens (group 17) have 7 valence electrons and tend to gain 1 electron to complete their octet.
\end{solucioneng}

%===================================================================
\newpage
\section{Inorganic Nomenclature}
\subsection*{Exam}
\examheader{Inorganic Nomenclature}
\setcounter{preg}{1}

\examexercise[1]{Name the following compounds using \textbf{both} the multiplier nomenclature and the oxidation-state nomenclature: \ch{Fe2O3}, \ch{CuH}, \ch{CaCl2}, \ch{MgO}.}
% CE: CE3.2

\examexercise[1]{Write the chemical formula of the following compounds: aluminium oxide, potassium hydride, magnesium chloride, sodium peroxide.}
% CE: CE3.2

\examexercise[2]{\ch{Na2O} and \ch{Na2O2} are two very different compounds even though they look similar. Explain, in your own words, what distinguishes an oxide from a peroxide, and explain why the subscript 2 in \ch{Na2O2} can never be simplified while the one in \ch{Ca2O2} can.}
% CE: CE1.1

\examexercise[1]{Explain the difference between \ch{HCl} and \ch{HCl(aq)}: what is each one called, and why does the name change even though it is ``the same compound''? Give a real-life example of where you might find \ch{HCl(aq)}.}
% CE: CE1.1

\examexercise[2]{Calcium hydroxide, \ch{Ca(OH)2}, is used by farmers in Castilla-La Mancha to correct the acidity of agricultural soil (a practice called \textit{encalado}). Knowing that the atomic masses are $\ch{Ca}=40$, $\ch{O}=16$, $\ch{H}=1$: \textbf{a)} calculate the molecular mass of \ch{Ca(OH)2}; \textbf{b)} calculate the mass percentage of oxygen in the compound.}
% CE: CE1.2

\examexercise[1]{In every binary compound, the element with a positive oxidation state is written on the left and the element with a negative oxidation state is written on the right. Explain, using the idea of electron transfer, why this rule makes chemical sense.}
% CE: CE3.2

\examexercise[1]{Classify each of the following formulas into its correct compound type (oxide, peroxide, hydride, hydroxide or binary salt): \ch{KOH}, \ch{FeCl3}, \ch{BaO2}, \ch{AlH3}, \ch{SO3}.}
% CE: CE1.1

\examexercise[1]{Quicklime, obtained by heating limestone in kilns (an industry historically present in several villages of Castilla-La Mancha), has the chemical formula \ch{CaO}. Write its name using the oxidation-state nomenclature and explain why no Roman numeral is needed. Then explain, in general terms, why this compound must be stored and handled carefully (hint: think about what happens when an oxide of this type comes into contact with water).}
% CE: CE1.3

\subsection*{Solutions}

\begin{solucioneng}
\textbf{Multiplier nomenclature / Oxidation-state nomenclature}
\begin{itemize}
    \item \ch{Fe2O3}: trioxide of diiron / iron (III) oxide
    \item \ch{CuH}: copper hydride (mono is omitted) / copper (I) hydride
    \item \ch{CaCl2}: calcium dichloride / calcium chloride (Ca has only one oxidation state, +2, so no Roman numeral is needed)
    \item \ch{MgO}: magnesium oxide (mono omitted) / magnesium oxide (Mg only has +2)
\end{itemize}
\end{solucioneng}

\begin{solucioneng}
\begin{itemize}
    \item Aluminium oxide: $\ch{Al}$ (+3), $\ch{O}$ (-2) $\rightarrow$ exchange oxidation states $\rightarrow$ \ch{Al2O3}
    \item Potassium hydride: $\ch{K}$ (+1), $\ch{H}$ (-1, metal hydride) $\rightarrow$ \ch{KH}
    \item Magnesium chloride: $\ch{Mg}$ (+2), $\ch{Cl}$ (-1) $\rightarrow$ \ch{MgCl2}
    \item Sodium peroxide: Na is group 1 $\rightarrow$ \ch{Na2O2} (never simplified)
\end{itemize}
\end{solucioneng}

\begin{solucioneng}
An \textbf{oxide} is formed by a metal (or nonmetal) and oxygen with oxygen in its usual -2 oxidation state; the subscripts can be simplified if there is a common factor (e.g. \ch{Ca2O2 -> CaO}). A \textbf{peroxide} contains the group \ch{O2^2-}, in which the two oxygen atoms are joined to each other by a special bond; that \ch{O2} unit is indivisible, so it can never be reduced to a single O, even if the metal subscript could otherwise be simplified. That is why \ch{Na2O2} keeps its \ch{O2} regardless of the fact that Na is also group 1, while \ch{Ca2O2} (two separate, independent oxide ions) simplifies to \ch{CaO}.
\end{solucioneng}

\begin{solucioneng}
\ch{HCl} (as a pure gas, not dissolved in water) is named \textbf{hydrogen chloride}, a nonmetal hydride. When this same compound is dissolved in water, forming an acidic aqueous solution, it is called \textbf{hydrochloric acid}, \ch{HCl(aq)}. The formula and the atoms are the same; what changes is the physical/chemical context (dissolved in water, behaving as an acid), so IUPAC nomenclature gives it a different name to communicate that difference. Real-life example: hydrochloric acid \ch{HCl(aq)} is naturally present in human gastric (stomach) juice, where it helps digest food.
\end{solucioneng}

\begin{solucioneng}
\textbf{a)} $M_w(\ch{Ca(OH)2}) = 40 + 2\times(16+1) = 40 + 34 = \qty{74}{\g\per\mole}$

\textbf{b)} There are 2 oxygen atoms in the formula, so:
\begin{equation*}
    \%\,\ch{O} = \dfrac{2 \times 16}{74} \times 100 = \dfrac{32}{74}\times 100 = 43.24\,\%
\end{equation*}
\end{solucioneng}

\begin{solucioneng}
When two elements combine, one of them loses electrons (it is oxidised, gets a positive charge) and the other gains those electrons (it is reduced, gets a negative charge). By convention, the element that loses electrons --- the one that becomes positive --- is written first, on the left, and the element that gains electrons --- the one that becomes negative --- is written on the right. This convention lets any chemist, just by reading the order of a formula, immediately know which element ``gave away'' electrons and which one ``received'' them.
\end{solucioneng}

\begin{solucioneng}
\begin{itemize}
    \item \ch{KOH}: hydroxide (metal + OH group)
    \item \ch{FeCl3}: binary salt (metal + nonmetal, not O, not H, not OH)
    \item \ch{BaO2}: peroxide (metal from group 2 + \ch{O2} unit)
    \item \ch{AlH3}: metallic hydride (metal + H)
    \item \ch{SO3}: oxide (nonmetal + oxygen)
\end{itemize}
\end{solucioneng}

\begin{solucioneng}
\ch{CaO} is named \textbf{calcium oxide}. No Roman numeral is needed because calcium only has one possible oxidation state, +2, so there is no ambiguity to resolve. Handling advice: oxides of this type react vigorously (and exothermically) with water to form the corresponding hydroxide (\ch{CaO + H2O -> Ca(OH)2}), releasing a lot of heat; because of this, quicklime must be stored away from moisture and handled with protective equipment, since contact with skin or eyes together with water can cause burns.
\end{solucioneng}

%===================================================================
\newpage
\section{Unit 7: Elements and Compounds}
\subsection*{Exam}
\examheader{Unit 7}
\setcounter{preg}{1}

\examexercise[1]{Sodium (\ch{Na}) and potassium (\ch{K}) belong to the same group of the periodic table. Explain, using the concept of valence electrons, why they show similar chemical properties.}
% CE: CE1.1

\examexercise[1]{Classify the type of chemical bond (ionic, covalent or metallic) present in each of the following substances, and justify each answer in one sentence: \ch{NaCl}, \ch{Cu} (solid metal), \ch{HCl}.}
% CE: CE1.1

\examexercise[1]{Explain, using an example, the difference between an \textbf{atom}, a \textbf{molecule} and a \textbf{mole}. Why was the mole created as a unit?}
% CE: CE3.2

\examexercise[2]{Calculate the molecular mass of aluminium sulfate, \ch{Al2(SO4)3}, knowing that $\ch{Al}=27$, $\ch{S}=32$, $\ch{O}=16$. Then calculate how many grams there are in \qty{0.5}{\mole} of this compound.}
% CE: CE1.2

\examexercise[1]{How many water molecules and how many hydrogen atoms are there in \qty{0.2}{\mole} of \ch{H2O}? Use Avogadro's constant, $N_A = \num{6.022e23}\,\unit{\per\mole}$.}
% CE: CE3.1

\examexercise[2]{We have \qty{2}{\mole} of an ideal gas at a temperature of \qty{300}{\kelvin} and a pressure of \qty{1.5}{\atmospher}. Calculate the volume occupied by the gas using the ideal gas law $pV = nRT$ (with $R = \qty{0.082}{\atmospher\litre\per\mole\per\kelvin}$).}
% CE: CE2.3

\examexercise[1]{What volume, in litres, do \qty{3}{\mole} of nitrogen gas (\ch{N2}) occupy under standard conditions (STP)? How many moles are contained in \qty{5.6}{\litre} of any ideal gas at STP?}
% CE: CE1.2

\examexercise[1]{A farming cooperative in Castilla-La Mancha uses ammonium sulfate, \ch{(NH4)2SO4}, as a nitrogen fertiliser. Knowing that $\ch{N}=14$, $\ch{H}=1$, $\ch{S}=32$, $\ch{O}=16$, calculate the mass percentage of nitrogen in this fertiliser and explain why this percentage is useful information for a farmer.}
% CE: CE1.3

\subsection*{Solutions}

\begin{solucioneng}
Elements in the same group of the periodic table have the same number of electrons in their valence (outermost) shell. Since chemical properties are determined mainly by valence electrons, sodium and potassium --- both having 1 valence electron (group 1) --- tend to react in the same way (e.g. both lose 1 electron to form a $+1$ ion), so they display similar chemical properties.
\end{solucioneng}

\begin{solucioneng}
\begin{itemize}
    \item \ch{NaCl}: \textbf{ionic bond} --- a metal (Na) and a nonmetal (Cl) combine; Na transfers an electron to Cl, forming \ch{Na+} and \ch{Cl-} ions held together by electrical attraction.
    \item \ch{Cu} (solid): \textbf{metallic bond} --- Cu atoms are all metallic; their valence electrons move freely among all the atoms forming an ``electron sea''.
    \item \ch{HCl}: \textbf{covalent bond} --- two nonmetals (H and Cl) share a pair of electrons to complete their valence shells.
\end{itemize}
\end{solucioneng}

\begin{solucioneng}
An \textbf{atom} is the smallest particle of a chemical element (e.g. one atom of oxygen, O). A \textbf{molecule} is a group of atoms chemically bonded together that behaves as a single unit and keeps the properties of the substance (e.g. one molecule of water, \ch{H2O}, made of 2 H atoms and 1 O atom). A \textbf{mole} is a unit of amount of substance: it is simply a very large ``bag'' that always contains \num{6.022e23} elementary entities (atoms, molecules, ions...). The mole was created because atoms and molecules are far too small and numerous to count or weigh individually in the laboratory, so chemists needed a practical unit that connects a countable number of particles with a mass we can actually measure in grams.
\end{solucioneng}

\begin{solucioneng}
\textbf{Molecular mass}:
\[
    \ch{Al2(SO4)3} \left\{\rule{0pt}{1.4cm}
        \begin{array}{lcl}
            \ch{Al} &=& 2 \times 27 = 54\,\text{uma} \\
            \ch{S} &=& 3 \times 32 = 96\,\text{uma} \\
            \ch{O} &=& 12 \times 16 = 192\,\text{uma} \\
            \cline{1-3}
            \multicolumn{3}{r}{\textbf{Total} = 342\,\text{uma}}
        \end{array}
    \right.
\]
Therefore $M_w(\ch{Al2(SO4)3}) = \qty{342}{\g\per\mole}$.

\textbf{Mass in \qty{0.5}{\mole}}: using $m = n \times M_w$:
\begin{equation*}
    m = 0.5\,\text{mol} \times 342\,\text{g/mol} = \qty{171}{\g}
\end{equation*}
\end{solucioneng}

\begin{solucioneng}
\textbf{Molecules}:
\begin{equation*}
    0.2\ \cancelred{\unit{\mole}\ \text{of \ch{H2O}}} \cdot \dfrac{\num{6.022e23}\ \text{molecules}}{1\ \cancelred{\unit{\mole}}} = \num{1.2044e23}\ \text{molecules of \ch{H2O}}
\end{equation*}
\textbf{Hydrogen atoms}: each \ch{H2O} molecule has 2 H atoms, so:
\begin{equation*}
    \num{1.2044e23} \times 2 = \num{2.4088e23}\ \text{atoms of H}
\end{equation*}
\end{solucioneng}

\begin{solucioneng}
Using $pV = nRT$, solving for $V$:
\begin{equation*}
    V = \dfrac{nRT}{p} = \dfrac{2 \times 0.082 \times 300}{1.5} = \dfrac{49.2}{1.5} = \qty{32.8}{\litre}
\end{equation*}
\end{solucioneng}

\begin{solucioneng}
Under STP, 1 mole of any gas occupies \qty{22.4}{\litre}.

\textbf{Volume of 3 mol}:
\begin{equation*}
    V = 3\ \cancelred{\unit{\mole}} \times \dfrac{\qty{22.4}{\litre}}{1\ \cancelred{\unit{\mole}}} = \qty{67.2}{\litre}
\end{equation*}
\textbf{Moles in \qty{5.6}{\litre}}:
\begin{equation*}
    n = \dfrac{5.6}{22.4} = \qty{0.25}{\mole}
\end{equation*}
\end{solucioneng}

\begin{solucioneng}
\textbf{Molecular mass of \ch{(NH4)2SO4}}:
\[
    \ch{(NH4)2SO4} \left\{\rule{0pt}{1.4cm}
        \begin{array}{lcl}
            \ch{N} &=& 2 \times 14 = 28\,\text{uma} \\
            \ch{H} &=& 8 \times 1 = 8\,\text{uma} \\
            \ch{S} &=& 1 \times 32 = 32\,\text{uma} \\
            \ch{O} &=& 4 \times 16 = 64\,\text{uma} \\
            \cline{1-3}
            \multicolumn{3}{r}{\textbf{Total} = 132\,\text{uma}}
        \end{array}
    \right.
\]
\textbf{Percentage of nitrogen}:
\begin{equation*}
    \%\,\ch{N} = \dfrac{28}{132}\times 100 = 21.2\,\%
\end{equation*}
This percentage tells the farmer exactly how much of the fertiliser's mass is actually the useful nutrient (nitrogen); it allows them to calculate the correct amount of fertiliser to apply per hectare to supply the crop with the nitrogen it needs, avoiding both under-fertilisation and the economic and environmental cost of over-fertilisation.
\end{solucioneng}

%===================================================================
\newpage
\section{Unit 8: Chemical Reactions}
\subsection*{Exam}
\examheader{Unit 8}
\setcounter{preg}{1}

\examexercise[1]{Propane (\ch{C3H8}) burns in the presence of oxygen (\ch{O2}) to give carbon dioxide (\ch{CO2}) and water (\ch{H2O}). Write the chemical equation for this reaction and balance it.}
% CE: CE3.2

\examexercise[1]{Explain, in your own words, what the Law of Conservation of Mass (Lavoisier's Principle) states, and give an everyday example (different from combustion) where you think mass is conserved even though it may not look like it at first sight.}
% CE: CE1.1

\examexercise[1]{Balance the following chemical equations: \textbf{a)} $\ch{Fe} + \ch{O2} \rightarrow \ch{Fe2O3}$; \textbf{b)} $\ch{Al} + \ch{HCl} \rightarrow \ch{AlCl3} + \ch{H2}$; \textbf{c)} $\ch{N2} + \ch{H2} \rightarrow \ch{NH3}$.}
% CE: CE3.2

\examexercise[2]{Magnesium reacts with oxygen to give magnesium oxide, according to the equation $2\,\ch{Mg} + \ch{O2} \rightarrow 2\,\ch{MgO}$. If \qty{60}{\g} of magnesium react with the necessary amount of oxygen, calculate the mass of magnesium oxide obtained ($\ch{Mg}=24$, $\ch{O}=16$).}
% CE: CE1.2

\examexercise[2]{Zinc reacts with sulphur to form zinc sulphide, \ch{ZnS}, according to the balanced equation $\ch{Zn} + \ch{S} \rightarrow \ch{ZnS}$. If \qty{6.5}{\g} of zinc react exactly with \qty{3.2}{\g} of sulphur: \textbf{a)} what mass of zinc sulphide is obtained? \textbf{b)} what is the mass percentage of zinc and of sulphur in zinc sulphide?}
% CE: CE1.2

\examexercise[1]{Explain what collision theory says about how a chemical reaction takes place, and mention the two conditions that a collision between particles must meet to be effective.}
% CE: CE2.3

\examexercise[1]{Natural gas, mainly methane (\ch{CH4}), is used for heating in many homes. It burns according to the equation $\ch{CH4} + 2\,\ch{O2} \rightarrow \ch{CO2} + 2\,\ch{H2O}$. Explain in words what this equation represents, and explain why it is important that this combustion takes place with enough oxygen (properly ventilated).}
% CE: CE1.3

\examexercise[1]{Calcium carbonate, \ch{CaCO3}, decomposes on heating (calcination) according to the equation $\ch{CaCO3} \rightarrow \ch{CaO} + \ch{CO2}$. If \qty{10}{\mole} of \ch{CaCO3} decompose completely, calculate: \textbf{a)} how many moles of \ch{CO2} are released; \textbf{b)} what volume, in litres, that \ch{CO2} occupies under standard conditions (STP), knowing that 1 mole of any gas occupies \qty{22.4}{\litre} at STP.}
% CE: CE3.1

\subsection*{Solutions}

\begin{solucioneng}
\begin{reaction*}
    C3H8 + O2 -> CO2 + H2O
\end{reaction*}
Balancing carbon, hydrogen and then oxygen:
\begin{reaction*}
    C3H8 + 5 O2 -> 3 CO2 + 4 H2O
\end{reaction*}
\end{solucioneng}

\begin{solucioneng}
The Law of Conservation of Mass states that, in a chemical reaction taking place in a closed container, the total mass of the reactants is always equal to the total mass of the products: matter is neither created nor destroyed, it is only transformed. Example: when an iron nail rusts, it seems to gain mass (rust weighs more than the original nail), but that extra mass actually comes from oxygen taken from the air that has combined with the iron; if we had weighed the closed system (nail + all the air/oxygen involved) before and after, the total mass would be the same.
\end{solucioneng}

\begin{solucioneng}
\begin{itemize}
    \item \textbf{a)} $4\,\ch{Fe} + 3\,\ch{O2} \rightarrow 2\,\ch{Fe2O3}$
    \item \textbf{b)} $2\,\ch{Al} + 6\,\ch{HCl} \rightarrow 2\,\ch{AlCl3} + 3\,\ch{H2}$
    \item \textbf{c)} $\ch{N2} + 3\,\ch{H2} \rightarrow 2\,\ch{NH3}$
\end{itemize}
\end{solucioneng}

\begin{solucioneng}
Since the exact (stoichiometric) amounts react, by the Law of Conservation of Mass:
\begin{equation*}
    m(\ch{Mg}) + m(\ch{O2}) = m(\ch{MgO})
\end{equation*}
We first calculate the moles of Mg and use the equation ratio $2\,\text{mol Mg} \rightarrow 2\,\text{mol MgO}$ (1:1 ratio):
\begin{align*}
    n(\ch{Mg}) &= \dfrac{60}{24} = 2.5\,\text{mol} \\
    n(\ch{MgO}) &= 2.5\,\text{mol} \quad (\text{ratio } 1{:}1) \\
    M_w(\ch{MgO}) &= 24 + 16 = 40\,\text{g/mol} \\
    m(\ch{MgO}) &= 2.5 \times 40 = \qty{100}{\g}
\end{align*}
We obtain \qty{100}{\g} of magnesium oxide.
\end{solucioneng}

\begin{solucioneng}
\textbf{a)} Since the masses given react exactly (in stoichiometric proportion), by conservation of mass:
\begin{equation*}
    m(\ch{ZnS}) = m(\ch{Zn}) + m(\ch{S}) = 6.5 + 3.2 = \qty{9.7}{\g}
\end{equation*}
\textbf{b)} Mass percentages:
\begin{align*}
    \%\,\ch{Zn} &= \dfrac{6.5}{9.7}\times 100 = 67.0\,\% \\
    \%\,\ch{S} &= \dfrac{3.2}{9.7}\times 100 = 33.0\,\%
\end{align*}
\end{solucioneng}

\begin{solucioneng}
Collision theory states that a chemical reaction occurs when the particles (atoms, ions or molecules) of the reactants collide with each other. For this collision to be \textbf{effective} --- that is, for it to actually break bonds and form the products --- two conditions must be met: \textbf{1)} the colliding particles must have enough energy, and \textbf{2)} they must collide with the correct (proper) orientation.
\end{solucioneng}

\begin{solucioneng}
The equation means that one molecule of methane reacts with two molecules of oxygen to produce one molecule of carbon dioxide and two molecules of water (or, in terms of moles: 1 mole of \ch{CH4} reacts with 2 moles of \ch{O2} to give 1 mole of \ch{CO2} and 2 moles of \ch{H2O}). It is important to have enough oxygen (good ventilation) because if the combustion happens without sufficient \ch{O2} (incomplete combustion), toxic carbon monoxide (\ch{CO}) can be produced instead of carbon dioxide, which is dangerous for people's health in a closed room.
\end{solucioneng}

\begin{solucioneng}
The balanced equation is $\ch{CaCO3} \rightarrow \ch{CaO} + \ch{CO2}$, a 1:1 mole ratio between \ch{CaCO3} and \ch{CO2}.

\textbf{a)} $n(\ch{CO2}) = n(\ch{CaCO3}) = \qty{10}{\mole}$ (1:1 ratio).

\textbf{b)} Under STP, 1 mol of any gas occupies \qty{22.4}{\litre}:
\begin{equation*}
    V(\ch{CO2}) = 10\ \cancelred{\unit{\mole}} \times \dfrac{\qty{22.4}{\litre}}{1\ \cancelred{\unit{\mole}}} = \qty{224}{\litre}
\end{equation*}
\end{solucioneng}

\end{document}
